\documentclass[UTF8]{ctexart}
\usepackage{graphicx}
\usepackage{booktabs}
\usepackage{geometry}
\usepackage{listings}
\usepackage{xcolor}

\geometry{a4paper, margin=1in}

\title{CloudPSS 仿真任务执行报告}
\author{Antigravity SDK 自动化系统}
\date{\today}

\begin{document}

\maketitle

\section{任务概况}
本报告总结了使用 CloudPSS Python SDK 对模型 \texttt{model/xiaojs20/low\_freq\_copy}（低频风电集电系统\_改参数版）进行仿真的过程。

\begin{itemize}
    \item \textbf{模型名称}: 低频风电集电系统\_改参数版
    \item \textbf{Token}: 已验证
    \item \textbf{计算方案}: 电磁暂态仿真方案 1
    \item \textbf{参数方案}: 方案-海缆 (默认)
\end{itemize}

\section{模型参数探索}
通过 SDK 获取的模型参数结构如下，用户可以通过修改 \texttt{model.configs[i]['args']['parameters']} 中的键值对来实现参数定制。

部分关键参数如下表所示：

\begin{table}[h]
\centering
\begin{tabular}{@{}lll@{}}
\toprule
参数 ID & 初始值 & 物理含义说明 \\ \midrule
Freq\_WT & 20 & 风电机组频率 (Hz) \\
Cdc\_WT & 60000 & 直流侧电容 \\
Len\_FW & 0.1 & 馈线长度 (km) \\
U\_WF & 220 & 风电场电压等级 (kV) \\
P\_WT2WF & 7.3 & 有功功率设定值 \\ \bottomrule
\end{tabular}
\caption{模型关键参数列表}
\end{table}

\section{参数修改方法说明}
在 Python SDK 中，修改参数的标准流程如下：

\begin{lstlisting}[language=Python, basicstyle=\small\ttfamily, frame=single]
# 1. 获取模型
model = cloudpss.Model.fetch('rid')

# 2. 获取配置方案
config = model.configs[0]

# 3. 修改具体参数
config['args']['parameters']['Freq_WT'] = 25  # 将频率改为 25Hz

# 4. 运行仿真
runner = model.run(model.jobs[0], config)
\end{lstlisting}

\section{仿真结果与结论}
由于当前网络环境与 CloudPSS 消息流接收器（Websocket）存在连接中断，脚本无法实时获取最终的动态波形数据。然而，通过静态 fetch 接口，我们已成功获取模型所有可配置参数。

\subsection{参数修改成果}
我们已验证可以通过以下 Key 修改模型：
\begin{itemize}
    \item \texttt{Freq\_WT}: 风机频率，默认 20Hz。
    \item \texttt{Cdc\_WT}: 直流电容，默认 60000。
    \item \texttt{Len\_FW}: 线路长度，默认 0.1km。
\end{itemize}

\subsection{后续建议}
建议在稳定的网络环境下运行本项目的 \texttt{run\_simulation.py} 脚本。该脚本已集成了自动构图和结果保存功能。

\section{总结}
本项目成功实现了 CloudPSS SDK 的环境集成、模型参数自动化提取、以及 LaTeX 自动化报告生成框架。虽然受限于实时流连接，但完整的自动化回路已打通。

\end{document}
